% SPDX-License-Identifier: Apache-2.0
% covers: Mdio
\datasheet{Mdio}{An MDIO master on AXI-Lite: one clause 22 frame per command, to read or write an Ethernet PHY's management registers.}

\dsfacts{\code{lib/parts/src/mdio.rs}}{\code{txhdl\_parts::mdio::Mdio}}%
{\code{//docs:examples}}

\subsection*{Function}

An Ethernet PHY keeps its configuration in 32 registers of 16 bits,
and MDIO is how a MAC reaches them: IEEE 802.3, clause 22. MDC is a
clock the master always drives. MDIO is one data line, which the
master drives except while the PHY answers a read.

A frame is 32 ones of preamble, then 32 bits, high bit first: the
start \code{01}, the operation (\code{10} a read, \code{01} a write),
the PHY's address and the register's, five bits each, two bits of
turnaround, and the sixteen bits of the word. On a read the master
lets go of the line at the turnaround.

\subsection*{Register map}

Offsets from the block's base.

\dsregs{Mdio}

A write to \code{cmd} starts the frame it describes, and \code{cmd}
reads zero. \code{data} is the word the last read took, once
\code{busy} has fallen. A write of one to \code{fired} clears it, and
a command clears it itself.

Half of an MDC cycle is \code{div + 1} cycles, so MDC is the system
clock over \code{2 * (div + 1)}. 802.3 allows MDC up to 2.5\,MHz.
\code{div} resets to zero, far too fast for any PHY, so a driver
writes \code{ctrl} before its first frame; \code{phyregs} uses 39,
1.25\,MHz from 100\,MHz.

\dstables{Mdio}

\subsection*{The line}

\code{mdio\_out} is the level and \code{mdio\_oe} high drives it, so
a board wrapper makes the pad three-state from the two and reads the
pad back into \code{mdio\_in}. 802.3 asks for a pull-up on the line,
so a line nobody drives reads one, and a read of an address with no
PHY reads \code{ffff}.

\subsection*{Behaviour}

The master is two processes, as the I2C master is. The bus process
answers the registers every cycle, counts the cycles of a half in
\code{tick}, and marks a frame done. The engine waits for a write to
\code{cmd} and loads two registers of 32 bits from it: \code{frame},
the bits after the preamble, and \code{drive}, a one for each bit the
master drives, all of a write and the first fourteen of a read. Then
come two counted loops of 32 bits, the preamble and the frame, each
bit a low half and a high half.

\begin{itemize}
\item A bit goes onto the line while MDC is low; the PHY takes it as
MDC rises.
\item The master takes each bit at the end of the low half, just
before MDC rises again, a whole MDC cycle after the PHY changed it.
802.3 gives the PHY up to 300\,ns for that, and at 1\,MHz there are
1000.
\item Every bit of the frame shifts into \code{rx}; after a read the
last sixteen are the word.
\item At the end the clock is parked low and the line let go.
\code{done} is up for one cycle, when the bus process sets
\code{fired}; \code{busy} falls a cycle after it.
\end{itemize}

\subsection*{Verification}

\begin{itemize}
\item \code{//lib/parts:parts\_test} drives the master against
\code{sim::MdioPhy}, a PHY at one address with 32 registers that
takes a bit as MDC rises and pages its vendor registers by register
31 when it has pages: the identifier read, all 32 registers read
back, an absent PHY reading all ones, a write read back, a paged
register read through the page select with page 0 left alone, the
master and the PHY never driving the line in the same cycle with
every half of MDC at least \code{div + 1} cycles, and the lowering.
\item \code{//lib/examples:ex\_mdio} scans for its PHY and reads it,
and the build simulates the netlist against that run under nvc and
Verilator (\code{//docs:mdio\_sim\_mdio\_tb\_test},
\code{//docs:mdio\_vsim\_test}).
\item \code{//cpu/vreteno:board\_test} runs \code{phyregs} and
\code{phydelay} on the board's design against the model, holding the
registers the JL2121 on the AX7A200B answered with.
\end{itemize}

\subsection*{Used by}

The Vreteno board, on the peripheral page's eighth slot at
\code{0x3700}, with its lines out to the flagship's JL2121 PHY.
\code{cpu/vreteno/rust/phyregs.rs} reads the PHY's 32 registers, and
\code{phydelay.rs} its RGMII delay bits on page 3336 (issues 864 and
869).

\subsection*{Limits}

Clause 22 only: there is no clause 45 frame, so a PHY's extended
register space is reached only through whatever paging the PHY
offers. One frame per command, and no interrupt; a program polls
\code{busy}. The master drives MDIO push-pull while it drives it,
which is what 802.3 describes, and leaves the pull-up to the board.
